\begin{document}
\lhead[ \leftmark ]{Informatik, Informatik \& Gesellschaft}
\rhead[Informatik]{Cyril Wendl, \today, Winterthur}

\title{Bewertungsraster Factsheet\\ \large Informatik \& Gesellschaft}

\maketitle

\section*{Hinweise zum Factsheet}

\textbf{Bewertungskriterien}:
\begin{itemize}
    \item (2 Pt) Inhaltlich korrekte und umfassende Zusammenfassung des Kapitels
    \item (1 Pt) Mindestens eine selbst erstellte \textbf{Illustration} (Bild, Diagramm, Grafik, usw.), die Ihr Factsheet sinnvoll ergänzt.
    \item (1 Pt) Formale Kriterien
    \begin{itemize}
        \item Abgabeformat: PDF
        \item Dateiname: \enquote{\texttt{[Klassenname] - K[Kapitelnummer] [Kapitelname] Factsheet.pdf}} (z.B. \enquote{\texttt{2d - K2 Daten Factsheet.pdf}})
        \item 3-4 A4-Seiten
        \item Korrekte Sprache, vollständige Sätze
        \item Schriftgrösse: 11 Punkte
        \item Vor- und Nachnamen aller Gruppenmitglieder im PDF enthalten
    \end{itemize}
    \item (1 Pt) Formulieren Sie am Ende Ihrer Zusammenfassung zwei interessante, reichhaltige Diskussions-Fragen (die nicht einfach mit ja / nein beantwortet werden können) oder provokativ-pointierte Aussagen, die zur Diskussion anregen. Versuchen Sie, Fragen zu formulieren, die Bezug zu Ihrer eigenen Lebenswelt (Interessen, Hobbies etc.) haben.
\end{itemize}


Das Fact Sheet enthält zudem spannende Fragestellungen, welche die Grundlage für die Gruppendiskussionen bilden

Geben Sie im Anhang (zählt nicht zu den 3-4 Seiten) an, wer für welche(n) Unterabschnitt(e) zuständig war.

\vspace{1em}

Das Factsheet ist spätestens eine Woche vor dem Vortrag abzugeben.


\newpage

\input{Grundlagen_Info/08_Gesellschaft/Feedback/Feedback_2d.tex}
\input{Grundlagen_Info/08_Gesellschaft/Feedback/Feedback_2e.tex}

\end{document}
